\scanpage{34}
Cela prouve (ii). Pour (i), on note qu'on sait déjà que si \(Z\) est alg. équiv. à 0, alors l'hom correspondant \(u:A_Y(k)\longrightarrow A^i_{\mathrm{alb}}(X)\) est nulle (Th. 1.6.), donc si \(Z\) est \(\tau\)-équivalent à 0, alors l'hom \(u\) est tel que \(mu=0\) pour quelque \(m>0\), donc \(u=0\) car \(A_Y(k)\) est \(m\)-divisible.

\uline{Corollaire} \unsure{2.7|2.2}. Pour qu'un cycle algébrique de codim \(i\) sur \(X\) soit alg. équiv. à 0, il f. et s. qu'il soit Albanese-équivalent à un diviseur de la forme \(Z(a_0)\), où \deleted{\(A\) est \unsure{une}} \(Z\) est un cycle de codim \(i\) sur un \deleted{VA \(A\),} produit \(A\times X\), \(A\) une variété abélienne, tel que \(\deleted{\unsure{\operatorname{Alb}}}a\leadsto Z(a):A\longrightarrow\deleted{\unsure{Z}}A^i_{\mathrm{alb}}(X)\) soit un hom de groupes.

\uline{Corollaire} 2.3. Pour tt entier \(n>0\), le groupe \(P^i(X)\) des \deleted{classes d'\unsure{équivalence}, à équiv. d'Alb.}\inserted{\(A^i_{\mathrm{alb}}(X)\), formé des classes de cycles} alg. équiv. à 0, est \(n\) \uline{divisible}.

\uline{Corollaire} 2.4. Soient \(Z,Z'\in A^*_{\mathrm{alb}}(X)\), dont l'un est alg. équiv. à 0 et l'autre \(\tau\)-équiv. à 0. Alors \(ZZ'=0\).
\mspage{33}
